\section{Data, metrics and paired statistics}
\label{sec:s:evaluation}

\paragraph{Dataset partitions.}
The neuronal references are public two-photon stacks used in DEEP$^2$ \cite{wijethilake2023,wijethilake2023data}. Raw sampling is $0.25\um$ laterally and $1\um$ axially. A $32\times128\times128$ crop is laterally area-resampled to $32\times96\times96$ and normalized to unit peak. Its planes are assigned the $0.32\um$ analysis spacing, giving $30.72\times30.72\times10.24\um^3$: lateral extent changes slightly and axial extent is compressed by a factor of 3.125. Reconstruction errors use these simulation coordinates. Test cells are SOM004 cells 1--2, SOM005 cell 1 and SOM006 cells 1--3: six cells from three animals. Each cell contributes two prespecified crops selected by projected foreground area times depth spread. After fixing reconstruction parameters, eight additional eligible crops per cell were sampled uniformly by location, excluding the two already selected. The resulting 48 crops were reconstructed without retuning. There are 60 test crops in the main comparison; other full-field studies use the 12 prespecified crops unless stated otherwise.

Denoisers and supervised optical/decoder models are trained on SOM009, a fourth animal, as specified in \cref{sec:s:reconstruction,sec:s:joint}. TV/PnP parameter selection uses two vascular crops from a different animal; supervised checkpoint selection uses 48 held-out neuronal volumes from SOM009. Conformal calibration uses a held-out stack of SOM009 and does not tune the reconstruction. No test volume enters optical training, denoiser training, scalar tuning or interval calibration.

\paragraph{Profile and centroid errors.}
Let $g_c=\sum_z x_{\mathrm{gt}}(c,z)$ and retain columns with $g_c>0.1\max_{c'}g_{c'}$. We set $\pi_c=g_c/\sum_{c'}g_{c'}$ over retained columns and normalize each nonnegative profile to unit mass. The one-dimensional W1 distance equals $\Delta z$ times the sum of absolute CDF differences; the main paper averages it with $\pi_c$. The implementation clips negative values to zero for the profile metric and divides by column mass plus $10^{-12}$. A zero-mass estimate therefore remains a zero profile rather than being assigned a fabricated depth distribution. The centroid error uses the same reference columns and weights, but the feed-forward evaluation computes centroids from raw network outputs; the W1 metric uses their nonnegative part. TV/PnP estimates are clipped before both metrics. For positive, normalized profiles, W1 is at least their absolute centroid difference; this inequality is not asserted for signed feed-forward centroids. A centroid can match when the profile does not, but similar reductions in the two errors do not establish a larger gain in profile shape.

PSNR uses the complete volume and SCICO's default reference range, $\max(\bx_{\mathrm{gt}})-\min(\bx_{\mathrm{gt}})$ \cite{balke2022scico}. Foreground SSIM uses a Gaussian-weighted 3D SSIM map ($\sigma=1.5$, intensity range one), averaged over reference voxels above 10\,\% of peak intensity after two binary-dilation steps. These intensity metrics complement W1, which normalizes each column and therefore does not penalize brightness errors by itself.

\paragraph{Depth-blind references.}
The flat estimate preserves each column's total brightness but distributes it uniformly over depth. The depth-blind oracle preserves that brightness and uses the reference crop's global axial profile for every column. The latter has access to the correct projected structure and global depth distribution but no local depth placement. On the 12 prespecified crops, their W1 errors are $2.266$ and $1.432\um$, with centroid errors $2.148$ and $1.322\um$, respectively. They are context for interpreting axial metrics, not reconstruction methods supplied with the actual measurements.

\paragraph{Pairing and uncertainty.}
For each setting, encoders share the specimen, direction subset, nominal dose and noise-generation protocol. We average the five draws within each crop. A relative change is the ratio of learned and comparator crop means minus one, not the mean of per-crop percentages. Confidence intervals use 10,000 cell-bootstrap resamples, keeping all crops of a sampled cell together. Cells from the same animal may also be dependent; six cells from three animals do not support a precise population-level guarantee. The intervals describe this finite evaluation and are not corrected for all exploratory comparisons.

\paragraph{Consistent reconstruction metrics.}
W1, centroid, PSNR and foreground SSIM in \cref{tab:s:e2e,tab:s:criteria60,tab:s:main} are computed from the same completed reconstruction rows. Each complete encoder/method configuration contains all 60 crops and five draws. The primary three-plate/four-method table contains 3,600 unique rows. The primary five-design/four-method table contains 6,000 unique rows, including the Fisher and coherence designs. Each criterion configuration has the same 60 crops and five draws. The second optical sampling seed adds 2,700 rows for our, Fisher and coherence plates under TV/PnP, using the same crops. Ancillary 12-crop studies are labeled separately. The reconstruction reruns use the recorded operators, parameters, directions and noise. Relative to earlier runs, the largest per-row differences are $0.019$\,dB/$0.0004\um$ for TV and $0.198$\,dB/$0.0204\um$ for PnP (PSNR/centroid); mean metrics agree at the reported precision. The all-column W1 evaluation used a first-row reproduction check for each configuration before completing its remaining solves. We report the resulting numerical estimates rather than assuming bitwise reproducibility across GPU runs.
